\section{An integral identity for the cubic permanent Hessian}

\subsection*{Auxiliary finding O-029: a compact independent certificate for F-028}
For a three-by-three matrix \(X=(x_{ij})\), write
\(P(X)=\operatorname{per}_3(X)\), \(D(X)=\det_3(X)\), and
\(\Pi(X)=\prod_{i,j}x_{ij}\).
The Hessian is taken in all nine entries. The following is
an identity of polynomials over \(\mathbb Z\):
\[
           \boxed{\quad
           \det\operatorname{Hess}_9 P(X)
                    =64\Pi(X)-2P(X)D(X)^2.
           \quad}                                             \tag{O-029.1}
\]
This identity is supplied as a short independently verified
algebraic certificate, without a historical priority claim.

\begin{proof}
Under row and column scalings
\(X\mapsto\operatorname{diag}(u_i)X\operatorname{diag}(v_j)\),
put \(\chi=\prod_i u_i\prod_j v_j\).
Both \(P\) and \(D\) scale by \(\chi\), while \(\Pi\)
scales by \(\chi^3\).
The entry-space linear transformation has determinant
\(\chi^3\).
Twice differentiating the identity \(P(gX)=\chi P(X)\)
therefore gives
\[
       (\det g)^2\det\operatorname{Hess}P(gX)
                    =\chi^9\det\operatorname{Hess}P(X).
\]
Thus the left side of (O-029.1) also scales by \(\chi^3\).
On the dense open set with nonzero first row and first
column, row and column scalings reduce to
\[
                       X=\begin{pmatrix}
                             1&1&1\\1&a&b\\1&c&d
                          \end{pmatrix}.
\]

For a row \(w=(w_1,w_2,w_3)\), put
\[
 M(w)=\begin{pmatrix}
       0&w_3&w_2\\w_3&0&w_1\\w_2&w_1&0
       \end{pmatrix}.
\]
If the three rows of \(X\) are \(u,v,w\), the Hessian has
the block form
\[
 \begin{pmatrix}
  0&M(w)&M(v)\\
  M(w)&0&M(u)\\
  M(v)&M(u)&0
 \end{pmatrix}.
\]
In the normalized matrix, \(u=(1,1,1)\),
\(M(u)=J-I\), and \(M(u)^{-1}=(J-2I)/2\).
The bottom six-by-six block has determinant \(-4\).
Its Schur complement consequently gives
\[
 \det\operatorname{Hess}P(X)=\frac12\det T,\qquad
 T=M(w)(J-2I)M(v)+M(v)(J-2I)M(w).
\]
For full reproducibility, its six upper-triangular entries are
\[
\begin{aligned}
 T_{11}&=-2(a-b)(c-d),&
 T_{22}&=-2(b-1)(d-1),&
 T_{33}&=-2(a-1)(c-1),\\
 T_{12}&=ad-a+bc+2bd+b-c+d,\\
 T_{13}&=2ac+ad+a+bc-b+c-d,\\
 T_{23}&=-ad+a-bc+b+c+d+2.
\end{aligned}
\]
Set \(e=ad+b+c\), \(o=bc+a+d\).
The three-by-three determinant formula
\[
 \det T=T_{11}T_{22}T_{33}+2T_{12}T_{13}T_{23}
             -T_{11}T_{23}^2-T_{22}T_{13}^2-T_{33}T_{12}^2
\]
with the displayed entries gives
\[
                          \det T=128abcd-4(e+o)(e-o)^2.
\]
Since \(P=e+o\), \(D=e-o\), and \(\Pi=abcd\) on this
normalized matrix, this proves (O-029.1) on the dense open
set. Both sides are polynomials with integer coefficients,
so equality there proves the identity over \(\mathbb Z\).
\end{proof}

In particular, a derivative can be checked without an
adjugate of a nine-by-nine Hessian:
\[
 \frac{\partial H_P}{\partial x_{ij}}
  =64\!\!\prod_{(k,l)\ne(i,j)}x_{kl}
      -2\,\frac{\partial P}{\partial x_{ij}}D^2
      -4PD\,\frac{\partial D}{\partial x_{ij}}.
                                                               \tag{O-029.2}
\]
The two derivatives on the right are the ordinary
two-by-two permanent and signed determinant cofactors.
Formula (O-029.2), at the pencil in F-028, independently
reproduces all 252 integer Jacobian entries and its
specified nonzero fourteen-by-fourteen minor.
The full polynomial identity, the compact derivative
calculation, and the coefficient counts in O-028/F-028
are checked in
\path{agents/obstructions/check_cubic_hessian_pair_identity.py},
with exact output
\path{agents/obstructions/cubic_hessian_pair_identity_certificate.json}.

\paragraph{Integral geometric meaning.}
Equation (O-029.1) is an equality of sections of the
same torus-character line over the integral affine
space of three-by-three matrices.
The normalization in the proof only inverts entry
coordinates on a dense open set over \(\mathbb Q\);
the resulting polynomial identity itself holds over
every commutative base ring.
This follows from equality of its integer coefficients,
and is stronger than asserting equality after one
characteristic-zero specialization.
The associated character grading is the
diagonalizable-group/comodule dictionary of SGA~3, I,
Propositions 4.7.2--4.7.3, pp.~34--36.
The identity alone does not imply the F-028 multiplicity
comparison: that further conclusion uses the nonzero
Jacobian minor, actual orbit realization, and
characteristic-zero highest-weight decomposition.
